\chapter{FUTURE ENHANCEMENTS}

\section{Multi-Crop Support and Adaptive Setpoints}
The designed hydroponics system is limited to a single crop which is lettuce (Lactuca sativa). At the present state of the system, the crop-specific setpoints such as the pH and EC ranges as well as the parameters including pump time period and fan runtime are only available for lettuce. A crop database with predefined parameter ranges, along with automatic profile switching based on the selected crop, can be implemented. In hydroponic cultivation, nutrient requirements can vary across different growth phases such as seedling, flowering, fruiting, vegetative, and harvest stages. So, an adaptive nutrient management system can be added based on the growth stage of the crop. 

\section{Integration of Machine Learning for Predictive Control}
The current system is a closed-loop system, resulting in the system being reactive to the deviation that has occurred. Machine learning enables predictive control, where the deviation is anticipated prior to the occurrence.  The key implementations could be nutrient consumption prediction, intelligent dosing, anomaly detection, and more. Using historical EC data, the system can predict how long it takes for nutrient levels to drop below the threshold, enabling dynamic dosing intervals. Instead of fixed dosing durations, the relationship between EC/pH change and pump activation time can be modelled to enable adaptive, non-linear dosing. When an anomaly such as gradual sensor drift occurs it can be detected easily or when sudden abnormal readings occur, they can be flagged as a fault in the system. 

\section{Improved Sensor Accuracy and Calibration Mechanisms}
The current system can only be manually calibrated since there is sensor drift with time. The improvements to be made in the future can be on automatic calibration, more precise industrial grade sensors, and redundancy of sensors to provide fault tolerance. This framework will self-check by using solutions known to have pH levels, such as pH 4, 7, and 10. It will test the pH of such solutions, and compare the measurements with the expected one and adjust itself, as needed. More than one sensor can also be used to measure the same parameter, such as two pH probes. In this manner, comparisons can be made between the outcomes of the two sensors and identify when one of them is not functioning properly. In case results are not consistent, the system will ignore such data. Lastly, by substituting cheap hobby sensors with high quality probes, it is possible to improve precision, stability, and dependability and make the system applicable to a long term and real life use.

\section{Mobile Application and User Interface Enhancement}
The system currently has a web dashboard which displays data and allows adjustment of threshold and dosing options. To make it more accessible and to increase interactivity, a mobile application may be added or substituted in place of the dashboard. The mobile application would enable real-time monitoring of key parameters such as pH, EC, temperature, and humidity through an intuitive and user-friendly interface. It would even send alerts to keep users informed immediately when parameters deviate from the predefined thresholds.  Besides monitoring, the application can also aid the operation of remote controls, where setpoints can be adjusted, crop profiles updated, or actuators, like pumps and fans, operated manually from anywhere. To sum up, the issue of the development of a mobile application greatly enhances the aspect of convenience, responsiveness and usability, rendering the system more applicable to real-life and large-scale implementations.

